\section{The least possible mean for a positive symmetric-power test}\label{sec:weights}
Write $t=2\cos\theta$ and
\[
 C_j(t)=U_j(t/2)=\frac{\sin((j+1)\theta)}{\sin\theta},
 \qquad d\mu(t)=\frac{\sqrt{4-t^2}}{2\pi}\,dt,\quad -2\le t\le2.
\]
At the endpoints use the continuous values. The $C_j$ are monic and
orthonormal for the Sato--Tate probability measure $\mu$. We use
$C_{-1}=0$, $C_0=1$, and $tC_j=C_{j+1}+C_{j-1}$.
Thus the coefficient of $C_0$ is the Haar mean $\int P\,d\mu$,
and the coefficient of $C_j$ is the multiplicity weight assigned to
the $j$th symmetric power. Since $C_j$ has parity $j$, the two
highest character coefficients of a polynomial are also its two
highest ordinary coefficients. The character product identity is
\[
 C_i C_j = \sum_{h=0}^{\min(i,j)} C_{i+j-2h}.
\]

\subsection{A quadrature certificate}
Fix $k\ge1$, let $u=(k+1)/k$, and define $Q_a=C_k+aC_{k-1}$
for $0\le a\le u$. Write its zeros as $z_1<\cdots<z_k$.
They are all interior to $[-2,2]$ except that $z_1=-2$ when
$a=u$. To see this for $k\ge2$, evaluate $Q_a$ at the zeros of
$C_{k-1}$, where $Q_a=C_k$ has alternating signs by the sine formula.
The values $Q_a(2)=k+1+ak$ and $(-1)^k Q_a(-2)=k+1-ak$ supply
the two outer roots, counting $-2$ when equality holds. The degree
gives simplicity. For $k=1$ the zero is $-a$.

For the cardinal interpolation polynomials $\ell_i$ at the $z_i$,
put $\lambda_i=\int \ell_i\,d\mu$. For $\deg V\le2k-2$ division
by $Q_a$ gives a quotient of degree at most $k-2$, orthogonal to
$Q_a$, and a remainder equal to $\sum_i V(z_i)\ell_i$. Therefore
\[
 \int V\,d\mu=\sum_i\lambda_i V(z_i),
 \qquad \lambda_i=\int\ell_i^2\,d\mu>0.
\]
The second identity follows by taking $V=\ell_i^2$.
This elementary form of quasi-Gaussian quadrature is what makes the
optimization certify its own lower bound.

\subsection{Odd degree}
\begin{proposition}\label{prop:oddfront}
Among monic polynomials of degree $2k-1$ nonnegative on $[-2,2]$,
the least mean is $(k+1)/k$. The unique minimizer is
\[
 R_k(t)=\frac{Q_u(t)^2}{t+2}.
\]
Every character coefficient of $R_k$ is strictly positive.
\end{proposition}
\begin{proof}
Since $Q_u(-2)=0$, $R_k$ is a monic polynomial, nonnegative everywhere
on the interval, with a simple endpoint zero and $k-1$ double
interior zeros. If $P$ is a competitor, $P-R_k$ has degree at most
$2k-2$, so the rule at $a=u$ gives
$\int(P-R_k)\,d\mu=\sum_i\lambda_i P(z_i)\ge0$.
Equality forces the endpoint zero and every interior double zero.
Together with monicity they determine $R_k$. Writing $[C_j]R_k$ for
its character coefficients, their values are
\begin{align*}
 [C_{2j}]R_k&=\frac{k(k+1)-j(j+1)}{k^2}
                  &&(0\le j\le k-1),\\
 [C_{2j-1}]R_k&=\frac{j^2}{k^2}
                  &&(1\le j\le k).
\end{align*}
For a direct check multiply these coefficients by $t+2$: the coefficient
at $j$ becomes $2r_j+r_{j-1}+r_{j+1}$, with missing terms set to zero.
This is $1+u^2$ at even $j<2k$, $2u$ at odd $j<2k$, and $1$
at $2k$, the coefficients of $Q_u^2$ by the product identity.
In particular $r_0=u$ and every displayed coefficient is positive.
\end{proof}

\subsection{Even degree with specified weight on the preceding power}
For even degree, minimizing the mean with only the top coefficient
specified yields $C_k^2$, which omits every odd symmetric power.
Prescribing the coefficient $b$ of $C_{2k-1}$ quantifies the cost of
including the preceding power. The other coefficients are free in the
minimization.

\begin{theorem}\label{thm:evenfront}
Let $P$ be monic of degree $2k$, nonnegative on $[-2,2]$, with
$[C_{2k-1}]P=b\ge0$. The least possible mean is
\[
 M_k(b)=
 \begin{cases}
  1+b^2/4,&0\le b\le 2u,\\
  1-u^2+ub,&b\ge 2u .
 \end{cases}
\]
The unique minimizing polynomial is
\[
 P_{k,b}=
 \begin{cases}
  Q_{b/2}^2,&0\le b\le2u,\\
  Q_u^2+(b-2u)R_k,&b\ge2u .
 \end{cases}
\]
For $b>0$ every character coefficient of $P_{k,b}$ is strictly
positive. Thus imposing this additional requirement preserves the
minimum and its attainment.
\end{theorem}
\begin{proof}
For $0\le a\le u$, subtract $Q_a^2+(b-2a)C_{k-1}Q_a$ from
$P$ to cancel the two highest terms. Applying the quadrature rule to
the difference and using orthonormality gives the exact identity
\[
 \int P\,d\mu=1-a^2+ab+\sum_{i=1}^{k}\lambda_i P(z_i).
\]
Maximizing the quadratic lower bound over $[0,u]$ selects
$a=\min(b/2,u)$ and gives $M_k(b)$. In the first range $Q_a^2$
attains equality. For $b\ge2u$ the displayed sum of $Q_u^2$ and
$(b-2u)R_k$ is nonnegative, has the prescribed top two coefficients,
and vanishes at all quadrature nodes, so also attains equality.
In the interior-node case equality forces $k$ double zeros and hence
the monic square. At $a=u$ equality forces divisibility by $R_k$.
The remaining monic linear factor is uniquely specified by the
coefficient $b$. This includes the common boundary $b=2u$.

The product identity gives $[C_{2k}]Q_a^2=1$, all lower even
coefficients $1+a^2$, and all odd coefficients $2a$.
Proposition~\ref{prop:oddfront} supplies the additional coefficients
in the second range.
\end{proof}

\noindent
For $0<b\le2u$ the same $M_k(b)$ is the sharp minimum when
\emph{every} odd coefficient is required to be at least $b$.
Indeed a competitor has top odd coefficient $b'\ge b$ and
$M_k(b')\ge M_k(b)$, while $Q_{b/2}^2$ has all odd coefficients
exactly $b$. For $b>2u$ Theorem~\ref{thm:evenfront} prescribes the
top odd coefficient, and its remaining odd coefficients are given by
the formulas above. Without a prescribed positive lower bound on any
odd coefficient, the infimum among strictly positive character
coefficients is $1$, approached as $b\downarrow0$. The unique
unconstrained minimizer $C_k^2$ shows this infimum is unattained
under strict positivity.

\paragraph{Relation to the sine-polynomial literature.}
Multiplication by $\sin\theta$ sends $P(2\cos\theta)=\sum c_j C_j(2\cos\theta)$
to the sine polynomial $\sum c_j\sin((j+1)\theta)$ nonnegative on
$[0,\pi]$. Hence a result about the top coefficient after normalization
by the first sine coefficient is directly relevant here. Rogosinski
and Szeg\H{o} introduced this extremal setting~\cite{rs}.
Dimitrov and Merlo~\cite{dm}, Theorems~1 and~2, give the individual
top-coefficient extremizers, including the odd-degree minimum above
and the even-degree example $C_k^2$. Their next-to-top coefficient
extremizer for odd sine degree is the mean-normalized case $b=2$ of
Theorem~\ref{thm:evenfront}. The joint formula fixes both highest
coefficients and describes the transition at $2+2/k$. Royster and Suffridge also determined the low-degree coefficient region for
typically real cubics, which under the same sine substitution covers
$k=1$ here.\cite{roysuff} We give the joint formula in all degrees and
make no assertion of bibliographic priority.
